\documentclass{article}
\usepackage[utf8]{inputenc}
\usepackage[T2A,T1]{fontenc}
\usepackage[ukrainian,english]{babel}
\usepackage{amsmath,amssymb,amsthm,mathtools}
\usepackage{tikz-cd}
\usepackage{hyperref}

% Minimal local journal macros (compatible with Groupoid Infinity vol4 style)
\usepackage{enumitem}
\usepackage{listings}
\usepackage{amsfonts}
\usepackage{tikz}

\theoremstyle{definition}
\newtheorem{definition}{Definition}
\newtheorem{theorem}{Theorem}
\newtheorem{lemma}{Lemma}
\newtheorem{proposition}{Proposition}
\newtheorem{remark}{Remark}
\newtheorem{example}{Example}
\newtheorem{corollary}{Corollary}
\newtheorem{conjecture}{Conjecture}

\newcommand{\RP}{\mathbb{RP}}
\newcommand{\CP}{\mathbb{CP}}
\newcommand{\R}{\mathbb{R}}
\newcommand{\C}{\mathbb{C}}
\newcommand{\Hom}{\mathrm{Hom}}
\newcommand{\id}{\mathrm{id}}
\newcommand{\PGL}{\mathrm{PGL}}
\newcommand{\GL}{\mathrm{GL}}
\newcommand{\SL}{\mathrm{SL}}
% \O is already defined (slashed O). Use \Ort for the structure group.
\newcommand{\Ort}{\mathrm{O}}
\newcommand{\Hilb}{\mathcal{H}}
\newcommand{\shape}{\int}

\newif\ifincludeTOC
\includeTOCtrue

\lstset{basicstyle=\footnotesize\ttfamily,breaklines=true}

\begin{document}

\title{Issue XLVIII: Projective Spaces and Quantum Field Information}
\author{Namdak Tonpa $^1$}
\date{$^1$ Groupoid Infinity / \foreignlanguage{ukrainian}{Інститут формальної математики} \\
       \small Modal Homotopy Type Theory \\
       \today}
\maketitle

\begin{abstract}
We record how real projective geometry supplies an ambient space in which
quantum field data can be formulated projectively, how geometry limits and
signature changes act on correlators, and how ordinary cohomology,
differential cohomology, and quantum cohomology are three different
linearisations --- of geometry, of anomalies, and of enumerative geometry ---
which must not be identified.  Spitz's axioms for projective quantum fields
are recalled, together with the Freed-style geometric axioms and homotopical
algebraic quantum field theory.  Mannheim's conformal gravity is included as
a locally scale-invariant classical parent: its infrared reduction recovers
Einstein and Maxwell, and its de~Sitter reading of a global potential fits
galactic rotation curves without a dark-matter halo.  The further claim that
quantizations of this parent bifurcate toward M-theory and loop quantum
gravity is marked as a proposal of the companion note, not as a theorem.
The discussion is then placed on the language lattice of Groupoid Infinity:
a grand unification, if it is to respect that lattice, is a diagram of
Cosmic Cube vertices and functors, not a single Lagrangian and not a single
vertex.
\\[0.5em]
{\bf Keywords}: projective geometry, conformal gravity, quantum field theory,
quantum cohomology, de~Sitter, Bach equation, Modal HoTT, anomalies,
language lattice.
\end{abstract}

\ifincludeTOC
  \tableofcontents
\fi

\section{Projectivity of quantum theory}

Pure states of a quantum system form a projective space \(\mathbb{P}(\Hilb)\),
not a linear Hilbert space.  Global phases are physically irrelevant;
transition probabilities are given by the Fubini--Study metric (or its
higher-dimensional analogues); unitary representations of symmetry groups
act projectively.  The obstruction to lifting a projective representation
to a linear one is cohomological and, in the modern geometric language,
is measured by an invertible field theory (an anomaly).

The same philosophy extends from quantum mechanics to quantum fields.
Homogeneous four-dimensional Klein geometries embed into the ambient real
projective geometry \((\RP^4,\PGL_5(\R))\).  Inside this ambient space
deformations, conjugacy limits and contractions become matrix operations;
coordinate singularities disappear and correlators remain well-defined.

\section{Projective quantum fields (Spitz axioms)}

\begin{definition}[Projective quantum field, after Spitz]
A \emph{projective quantum field} on the ambient Klein geometry
\((\RP^4,\PGL_5(\R))\) is a triple
\[
\bigl(\rho,\;\{[\hat{\mathcal{O}}_\alpha([x])]\}_{[x]\in\RP^4},\;
      (U_G)_G\bigr)
\]
consisting of
\begin{itemize}
  \item a finite-dimensional representation \(\rho\) of \(\PGL_5(\R)\),
        or of its complexification, governing tensor indices;
  \item an equivalence class of operator-valued tempered distributions on
        \(\RP^4\), defined only up to smooth nowhere-zero scalar prefactors
        \(C^\infty(\RP^4,\R_{\neq 0})\);
  \item for each homogeneous subgeometry
        \((\mathsf{X},G)\hookrightarrow(\RP^4,\PGL_5(\R))\), a projective
        unitary representation \(U_G\) of that structure group \(G\) on a
        Hilbert space \(\Hilb_G\).
\end{itemize}
The transformation law on each subgeometry is
\[
U_G([g])\,\hat{\mathcal{O}}_\alpha([x])\,U_G([g])^\dagger
  = \sum_\beta\rho_{\alpha\beta}([g^{-1}])\,
    \hat{\mathcal{O}}_\beta([g\cdot x]).
\]
We do \emph{not} require the \(U_G\) to assemble into a unitary representation
of the full non-compact group \(\PGL_5(\R)\).  That group acts on the base.
Its quantum implementation is required only along the isometry groups of the
subgeometries (Poincar\'e, de~Sitter, anti-de~Sitter, Galilei, Carroll, and
their contractions)~\cite{Spitz2024}.
\end{definition}

Restriction to \((\mathsf{X},G)\) recovers ordinary quantum fields on that
geometry.  Conjugacy limits of geometries --- on the Lie-algebra level,
In\"on\"u--Wigner contractions~\cite{InonuWigner} --- act on the base by the
ambient projective action.  Because the operators are equivalence classes of
tempered distributions, the limits of correlators stay in the same class of
objects.

\begin{proposition}[Geometry limits, after Spitz]
\label{prop:limits}
Let a family of homogeneous subgeometries be realized as a continuous
conjugacy limit inside \((\RP^4,\PGL_5(\R))\).  Along this limit the
projective correlators remain bounded as tempered distributions.  When the
space-time geometry degenerates, the limiting correlators are supported on
the lower-dimensional strata of the degeneration (boundaries, light cones,
and similar).  Ultraviolet and infrared limits are instances, and they lower
the dimension of the support.
\end{proposition}

\noindent
This is the form in which we use the regularity theorem of~\cite{Spitz2024}.
We do not reprove it, and we do not strengthen it to a claim that the
limiting objects are finite-rank tensors.  The proposition systematises the
ambient-space techniques long used in conformal field theory.  It is not, by
itself, a construction of a conformal field theory.

Signature phenomena fit the same picture.  The standard Minkowski metric of
signature \((-,+,+,+)\) (or its opposite) is recovered by choosing an appropriate
quadratic form on the five-dimensional vector space whose projectivisation is
\(\RP^4\).  Changing the signature or performing a Wick rotation corresponds to
a continuous path in the space of quadratic forms; the projective field operators
transform accordingly and the correlators remain well-defined along the path.
Extending the ambient dimension or changing the projective embedding produces
the familiar families of isomorphic or semi-isomorphic models.

\section{Conformal gravity}

Philip Mannheim's conformal gravity is the locally scale-invariant completion
of the same geometric picture~\cite{Mannheim2011}.  The unique pure-gravity
action invariant under
\[
g_{\mu\nu}(x)\mapsto e^{2\alpha(x)}g_{\mu\nu}(x)
\]
is the Weyl-squared action
\[
I_{\mathrm{grav}}
  = -\alpha_g\int d^4x\,\sqrt{-g}\,
    C_{\lambda\mu\nu\kappa}C^{\lambda\mu\nu\kappa},
\]
where \(C_{\lambda\mu\nu\kappa}\) is the Weyl tensor.  Variation produces the
fourth-order Bach equation
\[
4\alpha_g\,W^{\mu\nu}=T^{\mu\nu},
\]
with \(W^{\mu\nu}\) traceless and divergenceless.  The coupling \(\alpha_g\) is
dimensionless, so the theory is power-counting renormalizable.  Power-counting
is not a construction of a unitary renormalized quantum theory.  A generic
fourth-order Lagrangian has an Ostrogradsky instability.  Mannheim's proposal
is that unitarity is recovered in the \(\mathrm{PT}\)-symmetric sector, via
the \(\mathrm{PT}\) inner product rather than by Hermiticity of a second-order
Hamiltonian~\cite{Mannheim2011,BenderMannheim}.  That recovery is a research
program.  This note uses it as the analytic setting of the classical theory,
not as a completed quantization.

Three structural claims of the conformal theory matter for the projective
setting.

\begin{enumerate}[label=\arabic*.]
  \item \textbf{No a priori classical curvature.}
        In the pure classical theory the only geometry is flat Minkowski.
        All curvature is induced by coupling to a quantized matter source.
        Gravity is not quantized independently: the matter commutators fix
        the gravitational ones, and the constraint
        \(T^{\mu\nu}_{\mathrm{grav}}+T^{\mu\nu}_{\mathrm{matter}}=0\)
        forces the zero-point energy density and pressure of the gravity
        sector to cancel those of the matter sector identically.
  \item \textbf{Cosmological constant cancellation.}
        When conformal symmetry is spontaneously broken and masses are
        generated dynamically, the zero-point structure readjusts and cancels
        the cosmological-constant term induced by the mass gap.  There is no
        fundamental \(\Lambda\); a de~Sitter curvature
        \(R_{\mu\nu}=3H^2 g_{\mu\nu}\) appears only as a conformal factor
        \(\Omega(t)=e^{Ht}\) (or \(\cosh(Ht)/H\)) on a flat metric,
        \[
          ds^2=\Omega^2(t,\mathbf{x})\,\eta_{\mu\nu}\,dx^\mu dx^\nu.
        \]
  \item \textbf{Global galactic potential.}
        The exact static spherically symmetric vacuum solution has the form
        \[
          ds^2=-B(r)\,dt^2+\frac{dr^2}{B(r)}+r^2\,d\Omega_2,
        \]
        with \(B(r)\) containing the Schwarzschild term, a linear potential
        \(\gamma r\), and a global quadratic term \(-\kappa r^2\).  The constant
        \(\kappa\) is universal across galaxies.  The Mannheim--O'Brien fits
        give a representative value
        \(\kappa\approx 9.54\times 10^{-54}\,\mathrm{cm}^{-2}\);
        what the projective theory needs is the existence of one such cosmic
        constant, not this particular fit.  Circular velocities in the
        weak-field limit are
        \[
          v^2
            = \frac{\beta^{*}c^{2}}{r}
            + \frac{\gamma^{*}c^{2}\,r}{2}
            - \kappa c^{2} r^{2},
        \]
        where \(\beta^{*}\) and \(\gamma^{*}\) are system parameters and
        \(\kappa\) is not.  There is no dark-matter halo in this formula:
        the linear and quadratic terms are global conformal
        effects~\cite{Mannheim2011,TonpaConf}.
\end{enumerate}

The conformal group of Minkowski four-space is \(\mathrm{SO}(4,2)\)
(equivalently \(\mathrm{SO}(2,4)\)), of dimension \(15\).  Its standard
projective model is the group of transformations of \(\RP^5\) preserving the
Klein quadric of null lines in \(\R^{2,4}\).  It is not the structure group
of Spitz's ambient pair \((\RP^4,\PGL_5(\R))\), whose group has dimension
\(24\).  A literal inclusion \(\mathrm{SO}(4,2)\subset\PGL_5(\R)\) is not
asserted.  The two geometries meet by embedding and by conjugacy limit:
conformal correlators are projective correlators restricted to a conformal
subgeometry, and the flat, de~Sitter, and anti-de~Sitter limits are geometry
limits already controlled by Proposition~\ref{prop:limits}.  Conformal invariance is the
residual symmetry of the light cone after that embedding is chosen.

\subsection{Einstein and Maxwell as infrared instances}

Maxwell theory in four dimensions is itself conformally invariant: its
stress tensor is traceless and couples directly to the Bach tensor.  The
Einstein--Hilbert action is not conformally invariant.  It appears only after
the conformal symmetry is broken, as the infrared effective action obtained
by integrating out the conformal factor (the route of 't~Hooft recalled
in~\cite{Mannheim2011}).  In the Groupoid Infinity packaging both theories
are functors out of the conformal parent that preserve the differential-geometric
structure but forget the local scale symmetry~\cite{TonpaConf,tHooft}:
\begin{itemize}
  \item Einstein gravity is the broken-phase, low-curvature instance;
  \item Maxwell theory is the conformally invariant gauge instance in
        dimension four.
\end{itemize}
Neither functor is the parent.  Both discard local Weyl invariance, which is
the information loss of the infrared transport.

\subsection{de~Sitter geometry, M-theory, and loop quantum gravity}

Three readings of the parent theory have to be kept apart.

\begin{enumerate}[label=\arabic*.]
  \item \textbf{de~Sitter, classical, due to Mannheim.}
        Pure de~Sitter is not put in by hand.  It is the conformal-to-flat
        metric above.  The measured cosmic \(\kappa\) embeds galaxies in an
        effective de~Sitter background, and perturbations about that
        background stay inside the Bach equation~\cite{Mannheim2011}.
  \item \textbf{M-theory, a proposal.}
        In the equivariant super homotopy type theory of
        Sati--Schreiber~\cite{sati2022anyonic,schreibershulman12cohesive} the
        Weyl tensor is replaced by higher-categorical analogues (string
        gerbes, five-brane charges).  The companion note proposes that, after
        flux quantization, the Bach equation becomes the equations of
        eleven-dimensional supergravity~\cite{TonpaConf}.  Mannheim's paper
        does not prove this.  It is a bifurcation conjecture.
  \item \textbf{Loop quantum gravity, a proposal.}
        The companion note proposes that the Bach tensor supplies a
        fourth-order regulator for the Ashtekar connection, that spin networks
        are a discrete shadow of the Bach equation, and that the Immirzi
        parameter is a modulus of conformal symmetry breaking rather than a
        new force~\cite{TonpaConf}.  This is not a theorem of loop quantum
        gravity, nor of Mannheim.
\end{enumerate}

The modality dictionary is fixed in Section~\ref{sec:modal}; the one-line
version used by the companion note overloads the sharp modality and is
withdrawn.  What survives is: shape sees the conformal class, flat sees
discrete spectra, and the Bach tensor is an infinitesimal curvature, not a
sharp global section.

\section{Linearisation via cohomology}

Three linearisations appear in this note.  They are not the same functor.

\begin{enumerate}[label=\arabic*.]
  \item Ordinary cohomology (singular, de~Rham) linearises the non-linear
        geometry of a space.
  \item Differential cohomology linearises gauge fields with connections,
        and it is the home of the anomaly and of the Bach class: the Bach
        tensor represents metric variations modulo local Weyl rescalings.
  \item Quantum cohomology linearises enumerative geometry.  The quantum
        product on \(H^*(\CP^n)\), and more generally on the cohomology of a
        Fano or symplectic manifold, is defined by Gromov--Witten invariants.
        For projective spaces the ring is explicit; its Stokes data relate
        enumerative geometry to Bridgeland stability.  As a Frobenius manifold
        it carries a flat pencil of connections~\cite{Dubrovin,CoatesCorti}.
\end{enumerate}

\begin{proposition}[Only the shadows agree]
\label{prop:cospan}
Let \(X\) be a closed symplectic manifold, or a smooth projective variety.
\begin{enumerate}[label=\roman*.]
  \item Quantum cohomology is a deformation of the cup product.  Specialising
        every Novikov variable of positive curve class to zero is a ring map
        \(\sigma\colon QH^*(X)\to H^*(X)\).  For \(\CP^n\),
        \[
          QH^*(\CP^n)=\mathbb{Z}[H,q]/(H^{n+1}-q),
          \qquad \sigma(q)=0,
        \]
        so \(\sigma\) is surjective and not injective.  The two are isomorphic
        as modules on the same generators, and not as rings.
  \item Differential cohomology sits in the Hopkins--Singer square
        \[
        \begin{tikzcd}
          \widehat{H}{}^n(X) \ar[r] \ar[d] &
            \Omega^n_{\mathrm{cl}}(X) \ar[d] \\
          H^n(X;\mathbb{Z}) \ar[r] &
            H^n_{\mathrm{dR}}(X)
        \end{tikzcd}
        \]
        The underlying-class map \(\upsilon\colon\widehat{H}{}^n(X)\to H^n(X;\mathbb{Z})\)
        forgets curvature.  It is not an isomorphism, and \(\widehat{H}\) does
        not factor through the shape: a shape equivalence need not preserve it.
  \item The resulting diagram is a cospan, not a triangle of isomorphisms,
        \[
          \widehat{H}{}^*(X)
            \xrightarrow{\;\upsilon\;}
          H^*(X)
            \xleftarrow{\;\sigma\;}
          QH^*(X).
        \]
        This is the only canonical comparison.  The left arrow forgets
        curvature.  The right arrow forgets curve counts.  No natural
        isomorphism identifies the three, because a homotopy equivalence
        induces an isomorphism on \(H^*\), and need not induce one on
        \(\widehat{H}{}^*\) or on \(QH^*\).
\end{enumerate}
\end{proposition}

In cohesive type theory the same obstruction is modal.  Ordinary cohomology
is a mapping space out of the shape, \(\pi_0\mathrm{Map}(\shape X, HA[n])\).
Differential cohomology is a mapping space in the smooth site, and does not
factor through \(\shape\).  Quantum cohomology is not a mapping space into a
single spectrum at all: it depends on a symplectic form and on a count of
holomorphic curves.  An identification would be a term equating these three
function types.  The cospan is the term that exists.

\begin{remark}[The comparison that is worth making]
The Dubrovin connection
\[
  \nabla_z
    = d + z^{-1}(-\star_q)
\]
is a flat meromorphic connection on the trivial bundle with fibre \(H^*(X)\).
Iritani's \(\widehat{\Gamma}\)-class matches its flat sections with a lattice
coming from \(K\)-theory~\cite{Iritani}.  That integral local system lives on
the Frobenius manifold, not on \(X\).  Differential cohomology of that base
can receive the lattice.  The resulting map runs from quantum data to a
differential-cohomological object on the moduli of the quantum connection.
It does not produce an isomorphism \(\widehat{H}{}^*(X)\cong QH^*(X)\), and it
does not compute either theory from Spitz correlators.
\end{remark}

Quantum cohomology is the prototype of a complete, axiomatised linearisation
of a non-linear projective theory.  It is not a computation of the cohomology
of Spitz correlators.  Those correlators are not Gromov--Witten invariants,
and this note does not identify \(QH^*(\RP^4)\) with them.  The honest
operation on the correlators is poorer and still schematic: their expectation
values are projective analogues of Wightman functions~\cite{Wightman}, and the
associated graded of a filtration by scaling, or the cohomology of a complex
of operators, recovers ordinary classes.  The full projective correlators keep
the information the graded object forgets.

In higher-categorical terms the correlators are expected to take values in an
\(\infty\)-category of spectra, or of modules over an \(E_\infty\)-ring.
Ordinary cohomology is then the composite of smashing with an
Eilenberg--MacLane spectrum and taking homotopy groups.  That expectation is
a target for an internalisation.  It is not a theorem of this note.

\section{Axiomatic systems}

Several interlocking axiomatic frameworks exist.

\begin{enumerate}[label=\arabic*.]
  \item \textbf{Projective quantum fields (Spitz).}
        The most direct axiomatisation of the geometric picture above.
        It generalises a fragment of the Wightman axioms to the projective
        ambient geometry, keeps a Hilbert space (so that comparison with
        textbook QFT remains easy), and makes geometry limits intrinsic.

  \item \textbf{Conformal gravity (Mannheim).}
        Local Weyl invariance, the Weyl-squared action, the Bach equation,
        and the proposed cancellation of zero-point energy and of the induced
        cosmological constant.  Einstein and Maxwell are infrared instances
        in Mannheim's sense.  de~Sitter geometry, in the conformal-to-flat
        form, is also his.  M-theory and loop quantum gravity are proposals
        of~\cite{TonpaConf}, not axioms of this item~\cite{Mannheim2011}.

  \item \textbf{Atiyah--Segal / Freed-style geometric axioms.}
        A (possibly projective) functor from a bordism category (with extra
        structure: orientations, spin structures, principal bundles, \ldots)
        to a symmetric monoidal category (vector spaces, Hilbert spaces,
        categories, spectra, \ldots)~\cite{Atiyah1988}.  Projectivity appears
        as the need for invertible field theories that measure the failure of
        the functor to be strictly linear~\cite{FreedAnomalies,FreedHopkins}.

  \item \textbf{Algebraic / homotopical AQFT.}
        Haag--Kastler nets of algebras~\cite{HaagKastler}, upgraded to operads
        or \(\infty\)-operads so that Einstein causality becomes
        homotopy-coherent.  This is the setting closest to Modal HoTT, and it
        is still an analogy until a model structure is named.

  \item \textbf{Quantum cohomology / Frobenius manifolds.}
        The enumerative-geometric linearisation of the cohomology of projective
        varieties (and more generally of symplectic manifolds).  Fully
        axiomatised and computable for projective spaces and weighted projective
        spaces.
\end{enumerate}

\section{Modal HoTT and Groupoid Infinity}
\label{sec:modal}

The homotopy type of real projective space and the Klein geometry used above
are different objects.
\[
  \RP^\infty \simeq \mathrm{B}(\mathbb{Z}/2) = K(\mathbb{Z}/2,1),
\]
while \(\RP^n\) is the \(n\)-truncation (the \(n\)-skeleton) of that
classifying space, not the classifying space itself.  In particular
\(\RP^4\not\simeq\mathrm{B}(\mathbb{Z}/2)\).  The pair
\((\RP^4,\PGL_5(\R))\) is a homogeneous geometry.  Its homotopy type is only
the underlying \(\infty\)-groupoid of the base.  Keeping the two apart is
required before any internalisation.

Cohesive Modal HoTT distinguishes four operations, and the physics does not
get to reassign them~\cite{schreibershulman12cohesive}.
\begin{itemize}
  \item The \emph{shape} modality \(\shape\) sends a space to its path
        \(\infty\)-groupoid.  A conformal class is shape data: a local Weyl
        rescaling is a path, and \(\Lambda\)-cancellation and the cosmic
        \(\kappa\) are cohomology classes of the shape, not fields on the
        chart.
  \item The \emph{flat} modality \(\flat\) discretises.  Areas and volumes of
        a spin-network shadow, if the loop-quantum proposal is adopted, are
        \(\flat\)-modal.  Flat does not see the conformal factor.
  \item The \emph{sharp} modality \(\sharp\) codiscretises.  Maps into a sharp
        object ignore cohesion.  Sharp is not the home of \(\kappa\).
  \item In differential cohesion the \emph{infinitesimal shape} \(\Im\) sees
        jets.  The Bach tensor is the curvature of that infinitesimal
        structure.  Its vanishing is a statement about \(\Im\), not about
        \(\sharp\).
\end{itemize}

Quantum fields, if internalised, are dependent types with a coherent action
of the structure group of the subgeometry.  Cohomology theories are maps into
Eilenberg--MacLane spectra, or into the differential refinements of Jack.
The internal languages are the lattice of~\cite{GIArch}, not a single
calculus.  The assignment used below is:
\begin{center}
\begin{tabular}{lll}
\hline
Fragment & Language & What is interpreted \\
\hline
Strict projective syntax, \(\Pi\) only & Henk & signatures, no cohesion \\
Inductive classical analysis & Laurent & the infrared, \(\kappa\), \(v^2\) \\
Cubical paths, no scale & Anders & underlying \(\infty\)-groupoid \\
Shape, flat, sharp, \(\Im\) & Felix & conformal class, Bach curvature \\
Spectra, differential cohomology & Jack & anomalies, Bach class \\
Linear types, branes, braids & Urs & the M-theory proposal \\
Presentations, face maps & Dan & the spin-network proposal \\
Directed interval & Ulrik & the STT arrow into Urs \\
\hline
\end{tabular}
\end{center}
The companion
note~\cite{TonpaConf} records the conformal parent in the vol.~4 format; it
does not discharge the conjectures flagged above.

\begin{remark}[Signature and dimension]
Changing the signature of the ambient quadratic form, or extending the
projective dimension, produces a continuous family of models.  Isomorphism
is the strong relation, inside one vertex of the cube.  Semi-isomorphism
means a weak equivalence after localisation at shape equivalences: the same
\(\shape\)-type, different charts, possibly different signature.  That is the
precise form of the physicist's signature flip and dimensional extension.
It is not an isomorphism of the field theories, and it does not identify
their Bach tensors.
\end{remark}

\section{Toward a GUT: the language lattice}

The architecture of the Groupoid Infinity languages~\cite{GIArch} rejects the
conjecture that one formal system can be the internal language of all of
mathematics.  The stack is a lattice of three spines, not a tower and not a
single abstract syntax tree:
\begin{itemize}
  \item \textbf{Spine~I} (Henk--Frank--Christine--Laurent): constructive type
        theory and classical analysis;
  \item \textbf{Spine~II} (Henk--Per--Anders--Felix--Jack--Urs): super homotopy,
        differential cohesion, and quantum spectra;
  \item \textbf{Spine~III} (Dan--Mike--Ulrik): directed and simplicial
        \(\infty\)-category theory.
\end{itemize}
Cross-spine maps are functors (\textbf{Lift}, \textbf{Core}, \textbf{Modal},
\textbf{Diff}, \textbf{Quantum}), each with a definite information loss.
The same obstruction applies to a grand unified theory.  A GUT written as one
Lagrangian in one signature is the physical form of the monolithic AST: it
asks a single syntax to carry discrete spectra, directed higher category
theory, and cohesive geometry at once.  The projective and conformal material
of this note is one proposal for how to place that material on the lattice.
It is not a replacement parent theory.

\begin{definition}[GUT as a diagram of vertices]
A grand unification, in the sense compatible with~\cite{GIArch}, is not a
terminal syntax and not a single vertex of the Cosmic Cube.  The cube has
eight vertices; physics occupies several of them at once.  We use the vertex
table of~\cite{GIArch}, and we place the theories of this note as follows.
\begin{center}
\begin{tabular}{lll}
\hline
Theory & Vertex & Status \\
\hline
Wightman theory on a fixed background & \((\Delta,\mathbf{2},\mathbf{1})\) & axiomatic fragment \\
Cohesive conformal geometry, Bach & \((\nabla,\mathbb{N},\mathbf{1})\) & Mannheim, classical \\
Anomalies, differential cohomology & Jack, toward \((\nabla,\mathbb{N},\mathbb{H})\) & Freed \\
Linear quantum parent & \((\nabla,\mathbb{N},\mathbb{H})\) & Urs; M-theory is a proposal \\
Spin-network shadow & \((\nabla,\mathbf{2},\mathbb{H})\) & Dan; LQG is a proposal \\
Einstein, Maxwell, rotation curves & forgetful to Spine~I & infrared instances \\
\hline
\end{tabular}
\end{center}
Collapsing this table onto \((\nabla,\mathbb{N},\mathbb{H})\) would repeat, one
dimension up, the monolith the architecture paper rejects.  The Urs vertex is
where the linear and stable data live.  It is not where the Bach equation is
proved, and it is not where a galaxy curve is computed.
\end{definition}

The three axes still match the three usual, and unsuccessful, shortcuts.

\begin{enumerate}[label=\arabic*.]
  \item \textbf{Strictness} \(\Delta\leftrightarrow\nabla\).
        Wightman theory on a fixed background is strict.  Projective fields
        with geometry limits, and conformal geometry with local Weyl
        rescalings, are modal.  Moving along this axis replaces a background
        metric by the shape modality.  It does not, by itself, quantize the
        Bach equation.  The Bach tensor is an \(\Im\)-curvature
        (Section~\ref{sec:modal}), not a shape class.
  \item \textbf{Groupoidality} \(\mathbf{2}\leftrightarrow\mathbb{N}\).
        A set of field histories is \(1\)-categorical.  The space of pure
        states of a finite-dimensional system is a complex projective space
        \(\mathbb{CP}^{n-1}\), with colimit
        \(\mathbb{CP}^{\infty}\simeq K(\mathbb{Z},2)\), not
        \(\mathrm{B}(\mathbb{Z}/2)\).  Real projective space classifies real
        lines.  It is the ambient Klein geometry of the Spitz axioms, not the
        state space.  The space of geometries is an \(\infty\)-groupoid.  An anomaly is the failure of a projective
        representation, or of a bordism functor, to lift to a linear one.  That
        is a class in the Anderson dual of a bordism theory, in the sense
        of~\cite{FreedHopkins}.  It is not, on present evidence, a class in
        the syntactic cohomology of a programming language.
  \item \textbf{Stability} \(\mathbf{1}\leftrightarrow\mathbb{H}\).
        Einstein--Hilbert gravity is resource-insensitive.  The quantum
        completion Mannheim asks for would land in spectra (Jack) and in
        linear types (Urs).  Landing there is a conjecture.  Vacuum-energy
        cancellation is, in the classical theory, an identity of sources.  It
        is not yet a theorem of the stable homotopy category.
\end{enumerate}

The spines specialise these vertices.  Each functor loses something named.

\begin{itemize}
  \item \textbf{Spine~I, the infrared.}
        Laurent interprets the macroscopic regime: reals, measures, limits.
        Einstein and Maxwell are the images of the Bach equation after the
        conformal factor is integrated out~\cite{tHooft} and the directed data
        are forgotten.  The loss is local Weyl invariance.  Galactic rotation
        curves are statements in this language.  They do not require a
        dark-matter primitive.  They also do not constrain the Urs vertex.
  \item \textbf{Spine~II, cohesion before stability.}
        Felix has \(\shape\), \(\flat\), \(\sharp\), and \(\Im\), with the
        meanings fixed in Section~\ref{sec:modal}.  Jack passes to spectra and
        differential cohomology: this is the cohomological shadow of anomalies
        and of the Bach class, not of the raw Spitz correlators.  Urs adds
        configuration spaces, braids, and linear types.  The M-theory reading
        --- higher analogues of the Weyl tensor, then eleven-dimensional
        supergravity after flux quantization --- is a proposal about this
        vertex~\cite{TonpaConf}.  The \textbf{Diff} and \textbf{Quantum}
        functors lose, respectively, the non-stable cells and the non-linear
        correlators.
  \item \textbf{Spine~III, the discrete directed shadow.}
        Dan has face maps, boundaries, and presentations, and no dependent
        types.  The loop-quantum proposal places spin networks here, as a
        presentation, with the Immirzi parameter as a modulus of conformal
        symmetry breaking~\cite{TonpaConf}.  Ulrik's simplicial type theory
        supplies the directed interval.  The architecture paper's arrow
        \(\mathrm{STT}\to\mathrm{Urs}\) is why this shadow is not required to
        factor through Anders alone.  The \textbf{Core} functor
        \(\mathrm{Ulrik}\to\mathrm{Anders}\) discards direction and retains
        the maximal \(\infty\)-groupoid.  That groupoid is the space of
        invertible morphisms.  It is not the projective space of states.
        \(\mathbb{P}(\Hilb)\) is a further quotient, of a unit sphere by
        \(U(1)\), and the two must not be named as one object.
\end{itemize}

The spectral sequence of~\cite{GIArch},
\[
  E_2^{p,q}
    = H^p\bigl(\mathcal{I};\,\mathcal{H}^q\bigr)
    \;\Longrightarrow\;
    H^{p+q}(\mathcal{L}),
\]
computes the cohomology of the syntactic CW-complex of the language lattice.
It does not compute the scattering of a unified field theory.  The abutment
is a syntactic invariant.  The following dictionary is a conjecture about a
comparison map, not an identification.

\begin{conjecture}[Comparison of cohomologies]
There is a comparison, still to be constructed, from the anomaly theory of
the projective and conformal data to the syntactic cohomology of their
presentations in Jack and Urs, with the following values on low degrees.
\begin{itemize}
  \item \(H^0\) receives quantities conserved by rewriting and by the
        classical theory together: gauge charges, tracelessness of
        \(T^{\mu\nu}\), the projective equivalence class of a field operator.
  \item \(H^1\) is where a comparison would have to send anomalies: a
        projective phase, a gravitational anomaly, a failure of the conformal
        factor to decouple.  Equality with the syntactic \(H^1\) is not
        claimed.
  \item \(H^2\) is where a comparison would have to send coherence
        obstructions: a residual cosmological constant, the mismatch between
        a local halo and a global \(\kappa\), the Ostrogradsky mode.  The
        \(\mathrm{PT}\) inner product and the zero-point cancellation are
        proposed fillings of those cells.  They are not extra matter fields,
        and they are not yet \(2\)-cells in a proved chain complex.
\end{itemize}
\end{conjecture}

\begin{remark}[What is not being claimed]
This is not a spectrum-generating algebra for the Standard Model, and it
does not produce a unification scale.  The projective ambient geometry is the
common base.  Conformal gravity is the classical parent at
\((\nabla,\mathbb{N},\mathbf{1})\), not at every vertex.  A proposed GUT that
needs a new primitive for every regime, and a proposed GUT that occupies only
one vertex, are the same mistake in two directions.
\end{remark}

\section{Conclusion}

Projective geometry is an ambient Klein geometry for quantum fields, not a
classifying space of the field theory.  Spitz's axioms make geometry limits
and signature changes into operations on equivalence classes of tempered
distributions; the limiting correlators stay bounded and concentrate on the
degenerate strata.  Mannheim's conformal gravity is the locally
scale-invariant classical theory adapted to the light cone of that picture.
Its Bach equation is fourth-order and power-counting renormalizable;
\(\mathrm{PT}\)-unitarity is a program.  Einstein and Maxwell are infrared
instances.  A global quadratic potential accounts for the form of galactic
rotation curves without a dark-matter primitive.  de~Sitter geometry arises
as a conformal factor on flat space.

M-theory and loop quantum gravity are not theorems of this parent.  They are
proposals, placed respectively at the Urs vertex \((\nabla,\mathbb{N},\mathbb{H})\)
and at the Dan vertex \((\nabla,\mathbf{2},\mathbb{H})\).  Ordinary
cohomology, differential cohomology, and quantum cohomology form the cospan
of Proposition~\ref{prop:cospan}: both refinements map to ordinary cohomology,
and neither arrow is invertible.  A grand unification compatible with
the language lattice is the diagram of these vertices together with the
functors and the information each functor loses.  It is not one Lagrangian,
and it is not one corner of the cube.

\begin{thebibliography}{99}

\bibitem{Spitz2024}
Daniel Spitz,
\emph{Quantum fields on projective geometries},
J.\ Phys.\ A: Math.\ Theor.\ 58 (2025) 205203;
arXiv:2403.17996.

\bibitem{Mannheim2011}
Philip D.\ Mannheim,
\emph{Making the Case for Conformal Gravity},
Found.\ Phys.\ 42 (2012) 388--420;
arXiv:1101.2186.

\bibitem{TonpaConf}
Namdak Tonpa,
\emph{Conformal Gravity: de~Sitter Geometry, M-Theory, Loop Quantum Gravity,
Einstein and Maxwell Equations},
Groupoid Infinity, vol.~4;
\url{https://groupoid.space/books/vol4/conformal.pdf}.

\bibitem{FreedHopkins}
Daniel S.\ Freed and Michael J.\ Hopkins,
\emph{Reflection positivity and invertible topological phases},
Geom.\ Topol.\ 25 (2021) 1165--1330;
arXiv:1604.06527.

\bibitem{FreedAnomalies}
Daniel S.\ Freed,
\emph{Anomalies and invertible field theories},
arXiv:1404.7224.

\bibitem{CoatesCorti}
Tom Coates, Alessio Corti, Yuan-Pin Lee, Hsian-Hua Tseng,
\emph{The quantum orbifold cohomology of weighted projective spaces},
Acta Math.\ 202 (2009) 139--193.

\bibitem{Dubrovin}
Boris Dubrovin,
\emph{Geometry of 2D topological field theories},
in \emph{Integrable Systems and Quantum Groups},
Lecture Notes in Math.\ 1620, Springer, 1996.

\bibitem{GIArch}
Groupoid Infinity,
\emph{Architecture of Mathematics: Internal Languages, Categorical Models,
and the Spectral Lattice of Formal Systems},
2026.

\bibitem{Iritani}
Hiroshi Iritani,
\emph{An integral structure in quantum cohomology and mirror symmetry
for toric orbifolds},
Adv.\ Math.\ 222 (2009) 1016--1079.

\bibitem{Wightman}
R.\ F.\ Streater and A.\ S.\ Wightman,
\emph{PCT, Spin and Statistics, and All That},
W.\ A.\ Benjamin, 1964.

\bibitem{InonuWigner}
E.\ In\"on\"u and E.\ P.\ Wigner,
\emph{On the contraction of groups and their representations},
Proc.\ Natl.\ Acad.\ Sci.\ USA 39 (1953) 510--524.

\bibitem{HaagKastler}
R.\ Haag and D.\ Kastler,
\emph{An algebraic approach to quantum field theory},
J.\ Math.\ Phys.\ 5 (1964) 848--861.

\bibitem{Atiyah1988}
Michael Atiyah,
\emph{Topological quantum field theories},
Inst.\ Hautes \'Etudes Sci.\ Publ.\ Math.\ 68 (1988) 175--186.

\bibitem{tHooft}
Gerard 't~Hooft,
\emph{Probing the small distance structure of canonical quantum gravity
using the conformal group},
arXiv:1009.0669.

\bibitem{BenderMannheim}
Carl M.\ Bender and Philip D.\ Mannheim,
\emph{No-ghost theorem for the fourth-order derivative Pais--Uhlenbeck
oscillator model},
Phys.\ Rev.\ Lett.\ 100 (2008) 110402;
arXiv:0706.0207.

\bibitem{schreibershulman12cohesive}
Urs Schreiber and Michael Shulman,
\emph{Quantum gauge field theory in cohesive homotopy type theory},
in \emph{Proceedings of the Workshop on Quantum Physics and Logic}, 2012.

\bibitem{sati2022anyonic}
Hisham Sati and Urs Schreiber,
\emph{Anyonic topological order in twisted equivariant differential
K-theory},
arXiv:2206.13563.

\end{thebibliography}

\end{document}
